% Imports
\documentclass[12pt]{article}
\usepackage[margin=1in]{geometry}
\usepackage{
  amsmath, amsthm, amssymb
}

% Begin document
\setlength{\parindent}{0pt}
\begin{document}

% Title
\title{Econometrics II Workshop: Logistic Regression}
\author{Arturo Soberon}
\date{}
\maketitle

% Main
\setcounter{section}{0}

\section{The Setup: Modeling Default as a Bernoulli Outcome}

Suppose you work at a bank or fintech deciding whether to approve a loan
application. For each historical borrower $i$, we observe a binary outcome
$y_i \in \{0, 1\}$, where $y_i = 1$ means the borrower defaulted (didn't pay)
and $y_i = 0$ means they repaid on time.

If every applicant had the exact same fixed probability of default $p$, then
$y_i \sim \text{Bernoulli}(p)$, with probability mass function:
\begin{equation}
    P(Y_i = y_i) = p^{y_i}(1 - p)^{1 - y_i}
\end{equation}

\section{Conditioning on Borrower Characteristics}

Predicting a fixed default probability $p$ for everyone isn't useful---it
treats every applicant like a weighted coin flip with the exact same
probability $p$, ignoring who is actually applying. Instead, we observe a
$p \times 1$ vector of borrower characteristics $x_i$ (income, debt-to-income
ratio, credit bureau score, past delinquencies, plus an intercept) and want to
estimate each applicant's conditional probability of default:
\begin{equation}
    p_i = P(Y_i = 1 \mid x_i)
\end{equation}

\section{Why Not a Linear Probability Model?}

The most direct approach is to model default risk linearly: $p_i = x_i^\top \beta$.
The problem is that $x_i^\top \beta$ maps to the entire real line
$(-\infty, \infty)$, whereas a probability of default must live strictly inside
$[0, 1]$. Even ignoring heteroskedasticity, a model that tells your credit
committee an applicant has a $-15\%$ or $125\%$ chance of defaulting breaks
down as soon as you try to price risk or set an approval cutoff.

\section{The Sigmoid Function}

To keep our risk estimates well-behaved, we take our linear risk score
$z_i = x_i^\top \beta \in \mathbb{R}$ and pass it through a link function
that bounds the range of the output between $0$ and $1$, just the way a
probability should behave. In logistic regression, we use the \textbf{sigmoid}
(or logistic) function:
\begin{equation}
    \sigma(z) = \frac{1}{1 + e^{-z}}
\end{equation}
Setting $\hat{p}_i = \sigma(x_i^\top \beta)$ guarantees that every applicant's
predicted default probability lies strictly between $0$ and $1$.

\section{Sigmoid Function Properties}

This list is not exhaustive; it only covers the properties that will be useful
later in the lecture:
\begin{enumerate}
    \item \textbf{Range:} $\sigma: \mathbb{R} \to (0, 1)$, with
    $\lim_{z \to \infty} \sigma(z) = 1$ (near-certain default) and
    $\lim_{z \to -\infty} \sigma(z) = 0$ (near-zero default risk).
    \item \textbf{Symmetry:} $\sigma(z) = 1 - \sigma(-z)$, which cleanly links
    the probability of default to the probability of repayment.
    \item \textbf{Derivative:} $\sigma'(z) = \sigma(z)\bigl(1 - \sigma(z)\bigr)$.
\end{enumerate}

\textit{Quick check on Property 3:} Write $\sigma(z) = (1 + e^{-z})^{-1}$ and
apply the chain rule:
\begin{align}
    \sigma'(z) &= -(1 + e^{-z})^{-2}(-e^{-z}) \nonumber \\
               &= \frac{1}{1 + e^{-z}} \cdot \frac{e^{-z}}{1 + e^{-z}} \nonumber \\
               &= \sigma(z)\bigl(1 - \sigma(z)\bigr)
\end{align}

\section{Likelihood of the Loan Portfolio}

Given an i.i.d.\ historical dataset of $n$ loans $\{(x_i, y_i)\}_{i=1}^n$, the
joint likelihood of observing our portfolio's repayment outcomes as a function
of $\beta$ is the product of $n$ independent Bernoulli trials:
\begin{equation}
    \mathcal{L}(\beta) = \prod_{i=1}^n \hat{p}_i^{y_i} (1 - \hat{p}_i)^{1 - y_i},
    \quad \text{where } \hat{p}_i = \sigma(x_i^\top \beta)
\end{equation}

\section{Negative Log-Likelihood}

Taking the natural log turns the product across $n$ borrowers into a sum. As
you saw in Microeconomics, applying a strictly monotonic transformation ($\ln$)
preserves the argmax of an optimization problem---so maximizing
$\ln \mathcal{L}(\beta)$ gives the exact same $\hat{\beta}$ as maximizing
$\mathcal{L}(\beta)$. More importantly for computation, multiplying thousands
of probabilities in $(0, 1)$ produces a number so tiny that a computer rounds
it straight to zero (floating-point underflow).

Flipping the sign turns maximization into a standard minimization problem, and
dividing by $n$ gives our average negative log-likelihood (or binary
cross-entropy) objective function $J(\beta)$:
\begin{equation}
    J(\beta) = -\frac{1}{n} \ln \mathcal{L}(\beta)
             = -\frac{1}{n} \sum_{i=1}^n
               \Big[ y_i \ln(\hat{p}_i) + (1 - y_i) \ln(1 - \hat{p}_i) \Big]
\end{equation}

Think of $J(\beta)$ as a disutility or misery function that sets up the
computer's incentives: it gets penalized heavily whenever it assigns a low
default probability to someone who actually defaulted ($y_i = 1$), or a high
default probability to someone who paid on time ($y_i = 0$). To minimize its
misery, the algorithm is forced to push $\hat{p}_i$ toward $1$ for defaulters
and toward $0$ for reliable borrowers.

\section{Gradient Descent}

Unlike OLS, setting $\nabla_\beta J(\beta) = 0$ yields no closed-form
analytical solution because $\hat{p}_i$ is nonlinear in $\beta$. Fortunately,
$J(\beta)$ is globally convex---shaped like a smooth, multidimensional bowl
with a single lowest point---so we can reach the global minimum numerically
using \textbf{Gradient Descent}.

First, what is the gradient $\nabla_\beta J(\beta)$? It is simply the vector of
partial derivatives of our misery function with respect to each parameter:
\begin{equation}
    \nabla_\beta J(\beta) =
    \begin{bmatrix}
        \dfrac{\partial J}{\partial \beta_1} &
        \dfrac{\partial J}{\partial \beta_2} &
        \cdots &
        \dfrac{\partial J}{\partial \beta_p}
    \end{bmatrix}^\top
\end{equation}
Each entry tells you how much $J(\beta)$ changes if you nudge that specific
$\beta_j$ slightly upward. Put together into a single vector,
$\nabla_\beta J(\beta)$ always points in the direction of \textbf{steepest
ascent}---the direction in parameter space where our misery increases the
fastest.

Since we want to \emph{minimize} $J(\beta)$, we walk in the exact opposite
direction ($-\nabla_\beta J(\beta)$). Starting from an initial guess
$\beta^{(0)}$ and choosing a step size (learning rate) $\alpha > 0$, we
iteratively update our parameters:
\begin{equation}
    \beta^{(t+1)} = \beta^{(t)} - \alpha \nabla_\beta J(\beta^{(t)})
\end{equation}

Notice why this update rule gradually settles into the minimum:
\begin{itemize}
    \item \textbf{Sign check (1D intuition):} Suppose we only have one
    parameter $\beta_1$. If the slope $\frac{\partial J}{\partial \beta_1} > 0$,
    increasing $\beta_1$ makes things worse, so subtracting
    $\alpha \frac{\partial J}{\partial \beta_1}$ pushes $\beta_1$ to the left
    (downhill). If $\frac{\partial J}{\partial \beta_1} < 0$, subtracting a
    negative number pushes $\beta_1$ to the right (also downhill).
    \item \textbf{Natural braking:} When we are far from the bottom of the
    bowl, the slope is steep, so the gradient is large and we take big steps.
    As we approach the minimum, the surface flattens out and
    $\nabla_\beta J(\beta^{(t)}) \to 0$, meaning our steps automatically get
    smaller until the updates stop moving.
\end{itemize}

\section{Deriving the Gradient}

To execute that loop, we just need $\nabla_\beta J(\beta)$. Start with a single
borrower $i$, letting
\begin{equation}
    \ell_i(\beta) = -\Big[ y_i \ln(\hat{p}_i) + (1 - y_i) \ln(1 - \hat{p}_i) \Big]
\end{equation}
where $\hat{p}_i = \sigma(z_i)$ and $z_i = x_i^\top \beta$. Apply the chain rule:
\begin{equation}
    \frac{\partial \ell_i}{\partial \beta}
    = \frac{\partial \ell_i}{\partial \hat{p}_i}
      \cdot \frac{\partial \hat{p}_i}{\partial z_i}
      \cdot \frac{\partial z_i}{\partial \beta}
\end{equation}

Compute each piece individually:
\begin{align}
    \frac{\partial \ell_i}{\partial \hat{p}_i}
        &= -\left( \frac{y_i}{\hat{p}_i} - \frac{1 - y_i}{1 - \hat{p}_i} \right)
         = \frac{\hat{p}_i - y_i}{\hat{p}_i(1 - \hat{p}_i)} \\
    \frac{\partial \hat{p}_i}{\partial z_i}
        &= \sigma'(z_i) = \hat{p}_i(1 - \hat{p}_i) \\
    \frac{\partial z_i}{\partial \beta}
        &= x_i
\end{align}

Multiply them together---notice how the denominator cancels out cleanly:
\begin{equation}
    \frac{\partial \ell_i}{\partial \beta}
    = \frac{\hat{p}_i - y_i}{\hat{p}_i(1 - \hat{p}_i)}
      \cdot \hat{p}_i(1 - \hat{p}_i)
      \cdot x_i
    = (\hat{p}_i - y_i) x_i
\end{equation}

Average over all $n$ borrowers:
\begin{equation}
    \nabla_\beta J(\beta) = \frac{1}{n} \sum_{i=1}^n (\hat{p}_i - y_i) x_i
\end{equation}

Stacking the $x_i^\top$ rows into our $n \times p$ design matrix $X$ and
defining the $n \times 1$ vectors $\hat{p} = \sigma(X\beta)$ and $y$:
\begin{equation}
    \nabla_\beta J(\beta) = \frac{1}{n} X^\top (\hat{p} - y)
\end{equation}

Notice why this formula makes intuitive sense: in historical loan data, we
don't get to choose $X$ or $y$---those are fixed. The only lever we actually
control is $\beta$, which determines our predicted probabilities
$\hat{p} = \sigma(X\beta)$. The term $(\hat{p} - y)$ is simply a measure of the
prediction error we are making on each borrower, weighted by their observed
characteristics $X^\top$.

\section{Putting It Together}

Notice how familiar that gradient looks---it has the exact same structure as
OLS, except predicted default probabilities $\hat{p} = \sigma(X\beta)$ replace
$X\beta$. At each step $t$, all you need to do is evaluate
$\hat{p}^{(t)} = \sigma(X\beta^{(t)})$, compute the gradient
$\frac{1}{n} X^\top (\hat{p}^{(t)} - y)$, and plug it into the update rule to
get $\beta^{(t+1)}$ until convergence. Once trained, you score a new applicant
$x_{\text{new}}$ via $\hat{p}_{\text{new}} = \sigma(x_{\text{new}}^\top \hat{\beta})$
and approve the loan if $\hat{p}_{\text{new}}$ falls below your institution's
risk threshold.

% End document
\end{document}